\documentclass[10pt,a4paper]{amsart}
\usepackage{amsmath,amssymb,amsthm,graphicx}
\usepackage{verbatim}
\newenvironment{reference}{\comment}{\endcomment}
%\newenvironment{reference}{}{}
\newenvironment{slogan}{\comment}{\endcomment}
\newenvironment{history}{\comment}{\endcomment}

% For commutative diagrams we use Xy-pic
\usepackage[all]{xy}

% We use 2cell for 2-commutative diagrams.
\xyoption{2cell}
% Theorem environments.
%
\theoremstyle{plain}
\newtheorem{theorem}[subsection]{Theorem}
\newtheorem{proposition}[subsection]{Proposition}
\newtheorem{lemma}[subsection]{Lemma}

\theoremstyle{definition}
\newtheorem{definition}[subsection]{Definition}
\newtheorem{example}[subsection]{Example}
\newtheorem{exercise}[subsection]{Exercise}
\newtheorem{situation}[subsection]{Situation}

\theoremstyle{remark}
\newtheorem{remark}[subsection]{Remark}
\newtheorem{remarks}[subsection]{Remarks}

\numberwithin{equation}{subsection}

% Macros
%
\def\lim{\mathop{\mathrm{lim}}\nolimits}
\def\colim{\mathop{\mathrm{colim}}\nolimits}
\def\Spec{\mathop{\mathrm{Spec}}}
\def\Hom{\mathop{\mathrm{Hom}}\nolimits}
\def\Ext{\mathop{\mathrm{Ext}}\nolimits}
\def\SheafHom{\mathop{\mathcal{H}\!\mathit{om}}\nolimits}
\def\SheafExt{\mathop{\mathcal{E}\!\mathit{xt}}\nolimits}
\def\Sch{\mathit{Sch}}
\def\Mor{\mathop{\mathrm{Mor}}\nolimits}
\def\Ob{\mathop{\mathrm{Ob}}\nolimits}
\def\Sh{\mathop{\mathit{Sh}}\nolimits}
\def\NL{\mathop{N\!L}\nolimits}
\def\CH{\mathop{\mathrm{CH}}\nolimits}
\def\proetale{{pro\text{-}\acute{e}tale}}
\def\etale{{\acute{e}tale}}
\def\QCoh{\mathit{QCoh}}
\def\Ker{\mathop{\mathrm{Ker}}}
\def\Im{\mathop{\mathrm{Im}}}
\def\Coker{\mathop{\mathrm{Coker}}}
\def\Coim{\mathop{\mathrm{Coim}}}

% Boxtimes
%
\DeclareMathSymbol{\boxtimes}{\mathbin}{AMSa}{"02}

%
% Macros for moduli stacks/spaces
%
\def\QCohstack{\mathcal{QC}\!\mathit{oh}}
\def\Cohstack{\mathcal{C}\!\mathit{oh}}
\def\Spacesstack{\mathcal{S}\!\mathit{paces}}
\def\Quotfunctor{\mathrm{Quot}}
\def\Hilbfunctor{\mathrm{Hilb}}
\def\Curvesstack{\mathcal{C}\!\mathit{urves}}
\def\Polarizedstack{\mathcal{P}\!\mathit{olarized}}
\def\Complexesstack{\mathcal{C}\!\mathit{omplexes}}
% \Pic is the operator that assigns to X its picard group, usage \Pic(X)
% \Picardstack_{X/B} denotes the Picard stack of X over B
% \Picardfunctor_{X/B} denotes the Picard functor of X over B
\def\Pic{\mathop{\mathrm{Pic}}\nolimits}
\def\Picardstack{\mathcal{P}\!\mathit{ic}}
\def\Picardfunctor{\mathrm{Pic}}
\def\Deformationcategory{\mathcal{D}\!\mathit{ef}}
\usepackage{cleveref}
\crefname{originalreference}{original reference}{original references}
\Crefname{originalreference}{Original reference}{Original references}
\makeatletter
\newlabel{more-morphisms-lemma-quasi-finite-separated-pass-through-finite}{{05K0}{source}{}{sourcepointer-1}{}}
\newlabel{more-morphisms-lemma-quasi-finite-separated-pass-through-finite@cref}{{[originalreference][0][]05K0}{[1][0][]source}{}{}{}}
\makeatother
\crefname{originalreference}{original reference}{original references}
\Crefname{originalreference}{Original reference}{Original references}
\makeatletter
\newlabel{local-cohomology-lemma-finiteness-pushforwards-and-H1-local}{{0BK0}{source}{}{sourcepointer-2}{}}
\newlabel{local-cohomology-lemma-finiteness-pushforwards-and-H1-local@cref}{{[originalreference][0][]0BK0}{[1][0][]source}{}{}{}}
\makeatother
\crefname{originalreference}{original reference}{original references}
\Crefname{originalreference}{Original reference}{Original references}
\makeatletter
\newlabel{dualizing-lemma-depth}{{0AVZ}{source}{}{sourcepointer-3}{}}
\newlabel{dualizing-lemma-depth@cref}{{[originalreference][0][]0AVZ}{[1][0][]source}{}{}{}}
\makeatother
\title{Quasi-finite separated morphisms carrying a coherent module whose pushforward is coherent and acyclic are finite}
\author{The Stacks Project Authors}
\date{stacks-project:functors.tex:local-snapshot-20261012; GNU FDL 1.2 or later}
\begin{document}
\maketitle
\noindent\textit{Editorial note.} All delivered mathematical text is the author's own, copied verbatim from the pinned Stacks Project chapter file functors.tex; no statement or proof step was rewritten, reordered or removed.\par\medskip
\noindent\textit{Editorial note.} Span recomputation by label pairing (LANE\_BRIEF 10.2): in this chapter the author's proof environment follows the statement environment as a standalone block; the spec covers the complete environment pair (statement span, proof span) as recorded in the pool candidate.\par\medskip
\noindent\textit{Editorial note.} Occupation scan: combined claim\_key + statement\_anchor + generic-line scan over all registration files (active and withdrawn) found 0 hits for functors.tex (fu-occupied.json, 2026-10-14); consistent with the duty manager's six-check dispatch.\par\medskip
\ifdefined\hypertarget\hypertarget{sourcepointer-1}{}\fi
\noindent\textit{Source reference.} Original reference 05K0: sources/stacks-chow/tags.txt, source lines 7184--7184. Cross-chapter standard prerequisite cited by the author inside the counted span; it is used as a pointer and its own proof is not part of this item.\par\medskip
\ifdefined\hypertarget\hypertarget{sourcepointer-2}{}\fi
\noindent\textit{Source reference.} Original reference 0BK0: sources/stacks-chow/tags.txt, source lines 14543--14543. Cross-chapter standard prerequisite cited by the author inside the counted span; it is used as a pointer and its own proof is not part of this item.\par\medskip
\ifdefined\hypertarget\hypertarget{sourcepointer-3}{}\fi
\noindent\textit{Source reference.} Original reference 0AVZ: sources/stacks-chow/tags.txt, source lines 13702--13702. Cross-chapter standard prerequisite cited by the author inside the counted span; it is used as a pointer and its own proof is not part of this item.\par\medskip
\begin{lemma}
\label{lemma-characterize-finite}
Let $f : V \to X$ be a quasi-finite separated morphism of Noetherian
schemes. If there exists a coherent $\mathcal{O}_V$-module $\mathcal{K}$
whose support is $V$ such that $f_*\mathcal{K}$ is coherent and
$R^qf_*\mathcal{K} = 0$, then $f$ is finite.
\end{lemma}

\begin{proof}
By Zariski's main theorem we can find an open immersion
$j : V \to Y$ over $X$ with $\pi : Y \to X$ finite, see
More on Morphisms, Lemma
\ref{more-morphisms-lemma-quasi-finite-separated-pass-through-finite}.
Since $\pi$ is affine the functor $\pi_*$ is exact and faithful
on the category of coherent $\mathcal{O}_X$-modules.
Hence we see that $j_*\mathcal{K}$ is coherent and
that $R^qj_*\mathcal{K}$ is zero for $q > 0$.
In other words, we reduce to the case discussed in the next paragraph.

\medskip\noindent
Assume $f$ is an open immersion. We may replace $X$ by the
scheme theoretic closure of $V$. Assume $X \setminus V$ is nonempty
to get a contradiction. Choose a generic point $\xi \in X \setminus V$
of an irreducible component of $X \setminus V$. Looking at the situation
after base change by $\Spec(\mathcal{O}_{X, \xi}) \to X$ using flat base
change and using
Local Cohomology, Lemma
\ref{local-cohomology-lemma-finiteness-pushforwards-and-H1-local}
we reduce to the algebra problem discussed in the next paragraph.

\medskip\noindent
Let $(A, \mathfrak m)$ be a Noetherian local ring. Let $M$ be a finite
$A$-module whose support is $\Spec(A)$. Then $H^i_\mathfrak m(M) \not = 0$
for some $i$. This is true by
Dualizing Complexes, Lemma \ref{dualizing-lemma-depth}
and the fact that $M$ is not zero hence has finite depth.
\end{proof}

\end{document}
